\documentclass[output=paper]{langscibook}
\ChapterDOI{10.5281/zenodo.22880276}
\title{Introduction}
\author{Ulrike Mosel}
\abstract{}
\begin{document}
\maketitle\label{chap:25}
\begin{figure}[H]
    \includegraphics[trim=0 0 0 0,clip,width=.9\linewidth]{figures/25_1.jpg}
    \caption{Teop Island and the reef}
    \label{fig:25_1}
\end{figure}
The lagoon of the Teop region includes Teop Island and is very rich in fish and other marine animals.\footnote{For example, search the Semantic Domain of the \textsc{mted} for ``animals, fishes'' and the Part of Speech field for ``N2'' and find 131 entries. Or search the Semantic Domain for ``animals, shellfish'' and the Part of Speech field for ``N3'' and find 53 entries. These data were collected by the Teop team without any questionnaires.} Consequently, there has not been the need to develop special fishing techniques for the deep ocean. \textcite[185]{Shoffner:1976eg} remarks in his detailed chapter on Marine Resource Activities, ``The Teop are not great maritime travellers, seldom venturing more than a quarter of a mile into deep water.''

Decades ago traditional hand-made fishing nets and lines have been replaced with \textit{kave} \textit{sitiring} `nets made of plastic strings' and \textit{kara} \textit{sitiring} `plastic strings', but intensive team work with elderly fishing experts enabled the authors Enoch Magum, Naphtali Maion, Mark Tahi, and the graphic artist Neville Vitahi to write and illustrate the following chapters on catching turtles with the huge \textit{kave} \textit{goroto} `turtle net', see Chapters \ref{chap:26}, \ref{chap:27}, and \ref{chap:28}, the construction of the bamboo-frame net called \textit{kakavei}, see Chapter \ref{chap:29}, two net-fishing techniques called \textit{vari\-gomo} and \textit{seige}, see Chapters \ref{chap:30} and \ref{chap:31}, and angling with a kite to which the fishing line with a bait is attached, see Chapter \ref{chap:34}.\footnote{Compare Blackwood's descriptions of the ``net on frame'' and the ``fishing kite'' (\citeyear[342-346, 348-350]{Blackwood:1935im}). Shoffner reports that already during his fieldwork 1973-1974 turtle nets and fishing kites were no longer in use in the Teop area (\citeyear[208, 201]{Shoffner:1976eg}).}

Examples of fishing methods that were still practiced during our fieldwork are \textit{tanaa} `group fishing', see Chapter \ref{chap:32}, and \textit{vavaiku} `a kind of angling', see Chapter \ref{chap:33}, and those kinds of catching fish and gathering seafood that are only done by women on the reef during low tide. These are catching wrasses by hand, see Chapter \ref{chap:35}, fish poisoning, see Chapters \ref{chap:36}, \ref{chap:37}, and \ref{chap:38}, and digging sea-sausages, see Chapter \ref{chap:39}.

Further fishing methods like \textit{rooroto\textsuperscript{2}} `fish for tuna', \textit{taragao} `angling mackerels' and \textit{vavaakopo\textsuperscript{2}} `fishing with cole lamps', are described in the \textsc{mted}.\footnote{If you search the Semantic Domain for ``fishing'', you find 110 entries that are related to fishing methods. Fishing tools are listed in the Semantic Domain ``tools and weapons'' with references to 125 entries.}

The data for the texts and descriptions for the \textsc{mted} were collected by the Teop team without any questionnaires.
\end{document}